% ===========================================================================
% NOVA SUBSECÇÃO 2.2.3, a colocar a seguir a "ABORDAGENS TEMPORAIS E POR
% KEYFRAMES" e antes da secção "ATRIBUIÇÃO DE EQUIPA".
% ===========================================================================

\subsection{MÉTODOS APRESENTADOS AO DESAFIO SOCCERNET}
\label{sec:jnr_desafio}

As duas linhas anteriores descrevem a investigação publicada em conferências e revistas.
Existe, porém, um segundo corpo de trabalho que raramente chega a esses canais e que é
particularmente informativo para esta dissertação, o das soluções apresentadas às sucessivas
edições do desafio SoccerNet. Interessam aqui não pelas classificações que obtiveram, mas
pelas opções técnicas que revelam, e sobretudo pela mudança de paradigma que documentam ao
longo de três anos.

\subsubsection*{A família OCR em três fases}

A tarefa foi disputada isoladamente na edição de 2023, sobre um conjunto de \emph{tracklets}
já recortados, com 15 equipas e 157 submissões. Importa esclarecer que esses \emph{tracklets}
provêm do conjunto SoccerNet-Tracking e que a tarefa estava formulada como classificação, com
uma classe adicional para o número não visível, pelo que os valores dessa edição não são
comparáveis com os aqui reportados. O que dela se retira é o método.

Segundo os organizadores, a generalidade das equipas adotou um mesmo procedimento em três
fases. Primeiro uma deteção de texto sobre o recorte do jogador, para localizar o número na
imagem, com recurso a modelos como o DBNet++, o MMOCR, o DeepSolo ou o YOLO. Depois um
reconhecimento do texto localizado, por motores de OCR afinados para o efeito, tipicamente o
PP-OCRv3 ou o PaddleOCR. E por fim uma votação maioritária que consolida as leituras de todas
as imagens de um mesmo \emph{tracklet} numa decisão única. A solução vencedora seguiu esta
estrutura e concentrou o esforço nos dados, usando o DBNet++ para descartar do treino as
imagens sem número detetável, aplicando aumentação intensiva que incluía a fusão de várias
\emph{frames} numa só imagem, e combinando por votação cinco modelos de reconhecimento
distintos.

Duas observações se impõem. A primeira é que a agregação temporal por votação ao nível do
\emph{tracklet}, que também se adota neste trabalho, já era prática corrente da comunidade em
2023, e não constitui portanto uma novidade. A segunda é que toda esta família assenta em
tratar o número como texto arbitrário, exatamente a limitação que a Secção~\ref{sec:met_jnr}
procura ultrapassar.

\subsubsection*{Classificadores dedicados ao número}

Na primeira edição da reconstrução do estado de jogo, em 2024, as equipas melhor classificadas
abandonaram o OCR genérico e passaram a treinar modelos próprios, especializados no número de
camisola. Golovkin et al. (2025) organizam este panorama em duas categorias, a saber, as
adaptações de técnicas de OCR de uso geral, que localizam e depois leem o texto, e as soluções
desenhadas para o domínio desportivo. Sobre as primeiras notam que a exatidão se degrada com
jogadores distantes, com oclusões e com orientações laterais, e sobre as segundas que ganham
exatidão à custa de exigirem treino específico para cada conjunto de dados.

A própria solução vencedora ilustra bem a segunda categoria. Prescinde de detetar cada dígito
com uma caixa delimitadora e usa em seu lugar uma rede residual modificada, que recebe recortes
de 32 por 32 píxeis e possui duas cabeças de classificação, uma que decide se existe um dígito
inicial e qual, e outra que prevê o segundo dígito. A estrutura permite reconhecer o número
mesmo quando parte dele está ocluída, e dispensa a anotação de caixas por dígito. O
desequilíbrio entre números frequentes e raros é tratado com uma função de perda combinada.

A equipa classificada em segundo lugar seguiu caminho semelhante, treinando um detetor dedicado
de dígitos sobre dados públicos e cerca de mil imagens anotadas de propósito, e reconstruindo o
número a partir dos dígitos reconhecidos. Uma terceira equipa recorreu ao CLIP em duas fases,
com um modelo a decidir se há número visível e outro a identificá-lo, afinados sobre três mil
caixas anotadas manualmente.

O traço comum a todas é o custo de anotação. Cada uma destas soluções exigiu a construção de um
conjunto de treino próprio, o que representa um esforço incompatível com as condições em que
esta dissertação foi realizada.

\subsubsection*{A viragem para os modelos multimodais}

A edição de 2025 marca uma mudança clara. Os organizadores registam a aplicação bem-sucedida de
modelos de visão e linguagem à leitura do número, e as duas soluções mais relevantes nesse
sentido recorreram ao LLaMA-3.2-Vision e ao Qwen2-VL-Instruct.

A solução vencedora, descrita em detalhe por Oo et al. (2026), formula a identificação completa
do atleta como uma tarefa de geração autorregressiva multimodal. Em vez de treinar um
classificador por atributo, entrega ao modelo o recorte do jogador acompanhado de uma instrução
em linguagem natural que lhe pede a função, o número e a cor do equipamento, e lê a resposta
como uma sequência de símbolos. Os autores justificam a escolha por três razões, a saber, a
generalização a atributos nunca vistos, o que torna o reconhecimento aberto, a redução do
esforço de anotação, uma vez que basta ajustar a instrução, e a possibilidade de transpor a
solução para outros desportos.

Este mesmo trabalho fornece a comparação mais direta de que há registo entre as três estratégias,
por as avaliar dentro de um único \emph{pipeline} e sobre o mesmo conjunto de teste. A
combinação de modelos especializados por atributo, com o PRTreID para a função, o EasyOCR para o
número e uma rede residual para a cor, obtém 18.11\% de GS-HOTA. A substituição por um modelo de
visão e linguagem eleva o valor para 60.13\% com o CLIP e para 61.48\% com o LLaMA-3.2-Vision.
Ainda que estes valores digam respeito ao sistema completo, e não isoladamente à leitura do
número, a distância entre a primeira linha e as restantes é demasiado grande para ser atribuída
a outra coisa.

Uma segunda equipa aplicou o Qwen2-VL-Instruct de forma distinta, apresentando ao modelo uma
sequência longa de deteções dispostas em grelha e permitindo-lhe invocar uma função de ampliação
para inspecionar intervalos concretos antes de decidir. Em sentido contrário, uma das equipas
melhor classificadas manteve-se no reconhecimento de texto convencional, isolando primeiro a
região do tronco a partir dos pontos de interesse do ViTPose e aplicando depois o PARSeq, o que
mostra que a via clássica continua viável quando a região entregue ao leitor é a correta.

\subsubsection*{O número de camisola como estrangulamento reconhecido}

Vários participantes convergem no mesmo diagnóstico. A equipa da Universidade de Tóquio encerra
o seu relato reconhecendo que, tal como na solução de referência, a maior dificuldade permaneceu
no reconhecimento do número, e aponta os modelos multimodais como o caminho a seguir. Outra
equipa chega à mesma conclusão por via experimental, ao verificar que a calibração e a leitura do
número eram as etapas com maior margem no seu sistema.

A demonstração mais completa é, contudo, a de Oo et al. (2026), que decompõem a métrica oficial
por atributo dentro do seu próprio sistema. Avaliando apenas a posição no campo obtêm 79.52\%.
Acrescentando a função, o valor desce marginalmente para 79.20\%. Acrescentando em vez disso a
equipa, desce para 73.58\%. Acrescentando o número de camisola, cai para 64.70\%. Os autores
concluem que a leitura do número é o principal estrangulamento do sistema, e explicam-no pela
natureza da transmissão televisiva, onde os dígitos surgem desfocados, ocluídos ou visíveis em
poucas \emph{frames} de cada \emph{tracklet}.

Trata-se de confirmação externa e quantificada da premissa em que assenta esta dissertação, e que
motiva a escolha dos dois módulos a intervencionar.
